\section{Conclusion and Future Work}
This paper presented \revadd{a closed-loop framework} for dynamic robotic baggage handling in airport logistics environments. The proposed system tightly integrates \revadd{instance segmentation, spatio-temporal geometric modeling, motion planning, and active state regulation}. MSAM-Net improves segmentation robustness under severe tag interference through multi-scale aggregation and attention-guided mask refinement. PX3C and CAOS enhance 3D geometric reliability by suppressing residual outliers and stabilizing temporal OBB estimation. \revadd{The DKF-based motion planning module and relative-velocity planning strategy} enable stable grasp execution under conveyor motion, while active state regulation handles unstable or high-uncertainty configurations.

Experiments on a real robotic platform demonstrate that the proposed framework significantly improves precision-oriented grasping. While the general grasp success rate remains consistently high, \revadd{the precision success rate among successfully grasped suitcases improves from 62.1\% to 92.7\%.} \revadd{This improvement highlights the value of reliable geometric modeling and motion planning for precise placement in downstream baggage loading.} Although this work does not explicitly optimize packing policies, it provides the geometric and execution reliability required by high-density loading systems.

\revadd{The current validation is limited to predominantly rigid baggage on a structured conveyor at a nearly fixed speed. Performance with deformable baggage, severe clutter, variable conveyor motion, and different airport layouts remains to be established.}

\revadd{Future work will integrate tactile sensing to improve contact feedback and the handling of deformable baggage. We will also investigate packing planning that coordinates container layouts with grasping and placement to improve packing density in practical aircraft baggage loading.}
